\chapter{Elliptic classes and Euler degeneration}
\label{ed:chap:euler}

The elliptic characteristic class acts on cohomology by multiplication.
Where this multiplication is invertible, it transports geometric
correspondence operators by conjugation. At the specialization $y=1$, however, the class
becomes the top Chern class. On a connected positive-dimensional variety
with nonzero Euler characteristic, its multiplication has rank one.
To understand the specialization of the representation, we must therefore
retain the image of multiplication and its position inside cohomology.

We shall identify this image with an equivariant Euler ideal. Its quotient
is the cohomology of the projectivized tangent bundle, with the parameter
acting by the tautological class. This description turns the orders of
vanishing into lengths of nilpotent Jordan blocks.

Surfaces provide the
first example. The Hilbert square of a K3 surface exhibits several distinct
lengths. We then determine which operators preserve the image.
The usual Heisenberg vacuum relation obstructs an extension to the
unchanged ambient cohomology at $y=1$.

\section{The characteristic class and its Euler factor}
\label{ed:sec:factor}

Let $X$ be a smooth projective complex variety of complex dimension $d>0$.
We use rational singular cohomology, with its cup product and complex
orientation, and put
\[
 H_X=H^*(X;\mathbb Q),\qquad
 R_0=\mathbb Q[[\zeta]],\qquad R=\mathbb Q[[q,\zeta]],\qquad
 H_{X,R}=H_X\otimes_{\mathbb Q}R.
\]

The coefficient variables are even. We retain cohomological parity.
Multiplication by an inhomogeneous even class need not preserve the
integer cohomological grading. All power series use the indicated formal
topology. The variable $q$ records elliptic energy, and $y=e^\zeta$ is a
formal coordinate near $y=1$. We write
$\chi(X)=\sum_i(-1)^i\dim_{\mathbb Q}H^i(X;\mathbb Q)$
for the topological Euler characteristic.

For a vector bundle $V$, write
$\Lambda_tV=\sum_{j\geq0}t^j\Lambda^jV$ and
$\operatorname{Sym}_tV=\sum_{j\geq0}t^j\operatorname{Sym}^jV$.
The normalized elliptic tensor class is
\begin{equation}\label{ed:eq:elliptic-tensor}
 \mathcal E_{q,y}(X)=y^{-d/2}
 \bigotimes_{r\geq1}
 \left(
 \Lambda_{-yq^{r-1}}T_X^*\otimes
 \Lambda_{-y^{-1}q^r}T_X\otimes
 \operatorname{Sym}_{q^r}T_X^*\otimes
 \operatorname{Sym}_{q^r}T_X
 \right).
\end{equation}
The expression belongs to $K^0(X)_{\mathbb Q}[[q,\zeta]]$, where
$K^0(X)$ is the Grothendieck group of algebraic vector bundles.
The factor $y^{-d/2}$ means $e^{-d\zeta/2}$, also when $d$ is odd.
Its cohomological characteristic class and multiplication operator are
\begin{equation}\label{ed:eq:epsilon}
 \epsilon_X(q,e^\zeta)
 =\operatorname{td}(T_X)\operatorname{ch}(\mathcal E_{q,e^\zeta}(X)),
 \qquad E_X(a)=\epsilon_X(q,e^\zeta)\smile a.
\end{equation}

The Todd class supplies the Riemann--Roch normalization: integration of
$\epsilon_X$ gives the holomorphic Euler characteristic of the elliptic
tensor class. The operator $E_X$ retains more information than this
integral.

Let $x_1,\ldots,x_d$ denote formal Chern roots of $T_X$. Define
\begin{equation}\label{ed:eq:euler}
 e_\zeta(T_X)=\prod_{i=1}^d(x_i-\zeta)
 =\sum_{k=0}^d(-\zeta)^{d-k}c_k(T_X).
\end{equation}
This is the equivariant Euler class when fiber scaling has weight
$-\zeta$. We shall use the even power series
$s(w)=\sinh(w/2)/(w/2)$, whose constant term is one.
The factor at energy zero belonging to a single root is
\begin{equation}\label{ed:eq:single-root}
 e^{-\zeta/2}\frac{x(1-e^{\zeta-x})}{1-e^{-x}}
 =\frac{x\sinh((x-\zeta)/2)}{\sinh(x/2)}
 =(x-\zeta)\frac{s(x-\zeta)}{s(x)}.
\end{equation}
The remaining factors for that root form the series
\[
 \psi(q,\zeta,x)=
 \prod_{r\geq1}
 \frac{(1-e^\zeta q^r e^{-x})(1-e^{-\zeta}q^r e^x)}
      {(1-q^r e^{-x})(1-q^r e^x)}.
\]

\begin{theorem}\label{ed:thm:factor}
The elliptic class factors as
\begin{equation}\label{ed:eq:factor}
 \epsilon_X(q,e^\zeta)=e_\zeta(T_X)U_X(q,\zeta),
 \qquad U_X\in H_{X,R}^{\times},\qquad U_X(q,0)=1.
\end{equation}
The unit is given by
\begin{equation}\label{ed:eq:unit}
 U_X(q,\zeta)=\prod_{i=1}^d
       \frac{s(x_i-\zeta)}{s(x_i)}\psi(q,\zeta,x_i).
\end{equation}
In particular, the embedded ideals satisfy
\begin{equation}\label{ed:eq:ideal}
 \epsilon_XH_{X,R}=e_\zeta(T_X)H_{X,R}.
\end{equation}
The two multiplication complexes
$[H_{X,R}\xrightarrow{\epsilon_X}H_{X,R}]$ and
$[H_{X,R}\xrightarrow{e_\zeta}H_{X,R}]$ lie in resolution degrees
$-1,0$. Multiplication by $U_X$ on the source and the identity on the
target define an isomorphism from the first complex to the second.
\end{theorem}

\begin{proof}
For each power of $q$, only
finitely many factors contribute. Their denominators are $q$-adic units.
Positive cohomological degrees are nilpotent, so substitution of the
Chern roots into the series is defined. The factors involving $s$ are
units as well. At $\zeta=0$, every ratio in \eqref{ed:eq:unit} equals one.
The splitting principle and \eqref{ed:eq:single-root} prove
\eqref{ed:eq:factor}.

Multiplication by $U_X$ is an automorphism, which proves the equality of
ideals. The identity $e_\zeta U_X=\epsilon_X$ gives the chain map.
\end{proof}

Thus the ideal, its quotient, and the multiplication complex are
independent of elliptic energy after the displayed identification.
The conjugated correspondence operators will still depend on $q$.

\begin{example}[A surface with rationally trivial canonical class]
\label{ed:ex:surface}
Let $Y$ be a K3 or Enriques surface. Write $\omega\in H^4(Y;\mathbb Q)$
for the point class normalized by $\int_Y\omega=1$.
Since $c_1(T_Y)=0$ rationally and $c_2(T_Y)=\chi(Y)\omega$, we have
\[
 e_\zeta(T_Y)=\zeta^2+\chi(Y)\omega.
\]
Both surfaces have vanishing odd rational cohomology. On the span of
$1,\omega$, multiplication has matrix
\begin{equation}\label{ed:eq:surface-matrix}
 \begin{pmatrix}\zeta^2&0\\ \chi(Y)&\zeta^2\end{pmatrix},
\end{equation}
and on $H^2(Y;\mathbb Q)$ it is $\zeta^2$ times the identity.

The nonzero rational number $\chi(Y)$ is a unit of $R_0$.
Elementary row and column operations therefore give diagonal entries
$1,\zeta^4$ for \eqref{ed:eq:surface-matrix}, up to units. Consequently
\begin{equation}\label{ed:eq:surface-cokernel}
 \operatorname{coker}E_Y\cong
 R/(\zeta^4)\oplus\bigl(R/(\zeta^2)\bigr)^{b_2(Y)},
 \qquad b_2(Y)=
 \begin{cases}22,&Y\text{ a K3 surface},\\10,&Y\text{ an Enriques surface}.
 \end{cases}
\end{equation}
The length-four summand records the coupling of degree zero to top
degree. The rank-one specialization alone does not detect its length.
\end{example}

\section{The quotient as a projective bundle}
\label{ed:sec:projective}

Let $\pi:\mathbb P(T_X)\to X$ parametrize one-dimensional subspaces
of tangent fibers. Its tautological line is
$L=\mathcal O(-1)\subset\pi^*T_X$. Put $\lambda=c_1(L)$.
This convention determines the sign in the comparison with $\zeta$.

\begin{theorem}\label{ed:thm:projective}
There is a canonical isomorphism of $R$-superalgebras
\begin{equation}\label{ed:eq:projective}
 H_{X,R}/(\epsilon_X(q,e^\zeta))
 \ \cong\ H^*(\mathbb P(T_X);\mathbb Q)[[q]],
 \qquad a\longmapsto\pi^*a,\quad\zeta\longmapsto\lambda.
\end{equation}
On the right, $\zeta$ acts by cup product with the nilpotent class
$\lambda$.
\end{theorem}

\begin{proof}
The inclusion $L\hookrightarrow\pi^*T_X$ is a nowhere-zero section of
$\pi^*T_X\otimes L^*$. Hence its top Chern class vanishes:
\begin{equation}\label{ed:eq:projective-relation}
 \sum_{k=0}^d(-\lambda)^{d-k}\pi^*c_k(T_X)=0.
\end{equation}
The classes $1,\lambda,\ldots,\lambda^{d-1}$ restrict to a basis of
the cohomology of every fiber $\mathbb P^{d-1}$. Leray--Hirsch shows
that these classes form a basis over $H_X$. Since the relation is monic
up to sign, it gives the entire presentation
\[
 H^*(\mathbb P(T_X);\mathbb Q)
 \cong H_X[\zeta]/(e_\zeta(T_X)).
\]

In this graded presentation, $\zeta$ has degree two. The quotient has
finite top degree, so $\zeta$ is nilpotent. Completing in $\zeta$
therefore leaves the quotient unchanged. Extend in $q$ and apply
\eqref{ed:eq:ideal}.
\end{proof}

The same Euler polynomial occurs in zero-section self-intersection for
the tangent bundle with fiber weight $-\zeta$. Removing the zero section
and quotienting by scalar multiplication gives $\mathbb P(T_X)$.
Formula \eqref{ed:eq:projective} identifies the elliptic quotient with
the cohomology of this geometry. It does not require a choice of Chern
roots or a basis of cohomology.

\begin{example}[The sign for $\mathbb P^2$]\label{ed:ex:projective-plane}
Put $X=\mathbb P^2$ and $h=c_1(\mathcal O(1))$, so
$H_X=\mathbb Q[h]/(h^3)$ and $c(T_X)=1+3h+3h^2$. The relation is
\begin{equation}\label{ed:eq:p2-relation}
 \lambda^2-3h\lambda+3h^2=0.
\end{equation}

It can also be read directly on the complete flag variety. A tangent
line at a line $A\subset\mathbb C^3$ determines a plane $B$ containing
$A$, and is $\operatorname{Hom}(A,B/A)$. Thus, if
$a=c_1(A)$ and $b=c_1(B/A)$, then $h=-a$ and $\lambda=b-a$.

The flag relations include
$a^2+ab+b^2=0$ and $a^2b+ab^2=0$. Substitution gives
\eqref{ed:eq:p2-relation}. Changing the middle sign leaves
$6h\lambda$, which is nonzero since $1,\lambda$ form a basis over $H_X$.
This example detects the convention that is invisible
on surfaces with $c_1=0$.
\end{example}

\section{Elementary divisors and nilpotent lengths}
\label{ed:sec:smith}

Write $D=\dim_{\mathbb Q}H_X$. A matrix over the discrete valuation
ring $R_0$ with nonzero determinant admits invertible row and column
operations reducing it to
$\operatorname{diag}(\zeta^{\ell_1},\ldots,\zeta^{\ell_D})$,
where $0\leq\ell_1\leq\cdots\leq\ell_D$.
The integers $\ell_i$ are its Smith exponents. A zero exponent contributes
no cokernel. A positive exponent contributes $R_0/(\zeta^{\ell_i})$.

We use this theorem over $R_0$, then extend the resulting matrices to
$R$. We do not require a Smith theorem over a ring with two parameters.

\begin{proposition}\label{ed:prop:jordan}
Multiplication by $e_\zeta(T_X)$ on $H_X\otimes R_0$ is injective,
with determinant $(-\zeta)^{dD}$. Its positive Smith exponents are the
Jordan-block lengths of cup product by $\lambda$ on
$H^*(\mathbb P(T_X);\mathbb Q)$, and their sum is $dD$.
If $X$ is connected and $\chi(X)\ne0$, there is one zero Smith exponent
and $D-1$ positive exponents.
\end{proposition}

\begin{proof}
Order a homogeneous basis of $H_X$ by increasing degree. Every positive
Chern class raises degree, so the multiplication matrix is triangular
with diagonal $(-\zeta)^d$. Its determinant is nonzero in $R_0$.
The adjugate identity proves injectivity.

Each summand $R_0/(\zeta^\ell)$ is a single nilpotent Jordan block
for multiplication by $\zeta$. Theorem~\ref{ed:thm:projective}
identifies this endomorphism with cup product by $\lambda$.
Leray--Hirsch gives total dimension $dD$, which is also the valuation
of the determinant.

At $\zeta=0$ the matrix is multiplication by $c_d(T_X)$.
For connected $X$, let $\omega_X\in H^{2d}(X;\mathbb Q)$ satisfy
$\int_X\omega_X=1$. The operator vanishes on positive degrees
and sends $1$ to $\chi(X)\omega_X$. Its rank is one when
$\chi(X)\ne0$, so precisely one Smith exponent is zero.
\end{proof}

Let $C_X=\operatorname{coker}E_X$. Its derived fiber at $\zeta=0$
means $C_X\otimes_R^{\mathbf L}R/(\zeta)$, computed from a free
resolution. Theorem~\ref{ed:thm:factor} and the preceding proposition
give the resolution $[H_{X,R}\xrightarrow{\epsilon_X}H_{X,R}]$.
Thus a block of length $\ell>0$ has derived fiber
\[
 [\mathbb Q[[q]]\xrightarrow{0}\mathbb Q[[q]]]
\]
in resolution degrees $-1,0$. In the connected case with $\chi(X)\ne0$,
each of these degrees has rank $D-1$. The ordinary fiber and the derived
fiber have different information, but neither records the positive
lengths without retaining the parameter action.

For $X=\mathbb P^2$, the multiplication matrix in the basis $1,h,h^2$
is
\[
 \begin{pmatrix}
 \zeta^2&0&0\\-3\zeta&\zeta^2&0\\3&-3\zeta&\zeta^2
 \end{pmatrix}.
\]
The entries generate the unit ideal, the two-by-two minors generate
$(\zeta^2)$, and the determinant is $\zeta^6$. The exponents are
$0,2,4$. Equivalently, cup product by the tautological class on the
six-dimensional cohomology of the flag variety has blocks of lengths
two and four.

\section{The Hilbert square of a K3 surface}
\label{ed:sec:square}

Let $S$ be a projective K3 surface and $X=S^{[2]}$ its Hilbert scheme
of length-two subschemes. The degree shifts in the Hilbert-scheme
cohomology decomposition\cite{NakajimaHilbert1997} give, at length
two, $\operatorname{Sym}^2H^*(S)$ together with a second copy of
$H^*(S)$ shifted up by two. Therefore the even Betti numbers of $X$ are
\[
 (b_0,b_2,b_4,b_6,b_8)=(1,23,276,23,1),
\]
and its odd cohomology vanishes. The Hodge shift on the second summand
raises both indices by one, so
$h^{0,0}=h^{0,2}=h^{0,4}=1$ and $\chi(\mathcal O_X)=3$.

We use the Beauville--Bogomolov bilinear form $q_X$ on $H^2(X;\mathbb Q)$,
normalized by $\int_X\alpha^4=3q_X(\alpha,\alpha)^2$ and
$q_X(\alpha,\alpha)>0$ for an ample class $\alpha$.
The Chern--Fujiki identities give
\begin{equation}\label{ed:eq:chern-square}
 \int_Xc_4=324,\qquad
 \int_Xc_2\alpha\beta=30q_X(\alpha,\beta),\qquad
 \int_Xc_2^2=828.
\end{equation}
The first number also follows from the Betti numbers. The middle
identity is the length-two coefficient of the Chern--Fujiki generating
formula.\footnote{Y.~Cao, G.~Oberdieck and Y.~Toda,
\emph{Gopakumar--Vafa type invariants of holomorphic symplectic
4-folds}, Theorems~4.1--4.2,
\url{https://doi.org/10.1007/s00220-023-04882-8}.}

For the last number, Hirzebruch--Riemann--Roch and $c_1=0$ give
$720\chi(\mathcal O_X)=3\int_Xc_2^2-\int_Xc_4$.
Nondegeneracy of $q_X$ implies that multiplication by $c_2$ is an
isomorphism $H^2(X)\to H^6(X)$.

\begin{theorem}\label{ed:thm:square}
For $X=S^{[2]}$, the Smith exponents, with multiplicities, are
\begin{equation}\label{ed:eq:square-spectrum}
 0^1,\qquad 2^{23},\qquad 4^{276},\qquad 6^{23},\qquad 8^1.
\end{equation}
At every elliptic energy,
\[
 C_X\cong
 (R/(\zeta^2))^{23}\oplus(R/(\zeta^4))^{276}
 \oplus(R/(\zeta^6))^{23}\oplus R/(\zeta^8).
\]
Its derived fiber has rank $323$ over $\mathbb Q[[q]]$ in each of
resolution degrees $-1,0$.
\end{theorem}

\begin{proof}
The symplectic form identifies $T_X$ with $T_X^*$, so $c_1=c_3=0$
over $\mathbb Q$. Hence
$e_\zeta(T_X)=\zeta^4+\zeta^2c_2+c_4$.

Choose $\omega\in H^8(X;\mathbb Q)$ with $\int_X\omega=1$.
On the invariant span of $1,c_2,\omega$, the matrix is
\begin{equation}\label{ed:eq:square-core}
 \begin{pmatrix}
 \zeta^4&0&0\\
 \zeta^2&\zeta^4&0\\
 324&828\zeta^2&\zeta^4
 \end{pmatrix}.
\end{equation}
Its entries generate the unit ideal. All two-by-two minors are divisible
by $\zeta^4$, and the last two rows and first two columns have minor
$504\zeta^4$. Its determinant is $\zeta^{12}$. The exponents of this
block are consequently $0,4,8$.

Choose a basis $a_1,\ldots,a_{23}$ of $H^2(X)$. The classes
$c_2a_i$ form a basis of $H^6(X)$. Each pair $a_i,c_2a_i$ gives
\[
 \begin{pmatrix}\zeta^4&0\\\zeta^2&\zeta^4\end{pmatrix},
\]
whose exponents are $2,6$.

Finally, the kernel of $\int_Xc_2(-)$ on $H^4(X)$ has dimension
$275$ and is complementary to $\mathbb Qc_2$, because
$\int_Xc_2^2\ne0$. Multiplication on each of these directions is
$\zeta^4$. These invariant subspaces exhaust cohomology and yield
\eqref{ed:eq:square-spectrum}. The $323$ positive exponents sum to
$1296=4\cdot324$. The free resolutions of the positive blocks give
the derived fiber, and \eqref{ed:eq:factor} extends the decomposition
from $R_0$ to $R$.
\end{proof}

There is a geometric description of the same lengths. Write $C_i$ for
cup product by $c_i(T_X)$ on $H_X$. In the projective-bundle basis
$1,\lambda,\lambda^2,\lambda^3$, cup product by $\lambda$ has matrix
\begin{equation}\label{ed:eq:companion}
 \begin{pmatrix}
 0&0&0&-C_4\\
 I&0&0&0\\
 0&I&0&-C_2\\
 0&0&I&0
 \end{pmatrix}.
\end{equation}
The last column is the relation
$\lambda^4+c_2\lambda^2+c_4=0$.

The three-dimensional block \eqref{ed:eq:square-core} contributes
Jordan lengths $4,8$. Each pair from $H^2$ and $H^6$ contributes $2,6$.
Each remaining middle-degree direction contributes $4$.

For higher Hilbert schemes, the same projective-bundle theorem applies,
but the lengths require the corresponding Chern multiplication operators.
The length-zero Hilbert scheme is a point: its elliptic class is one,
its cokernel is zero, and its rank-zero tangent projectivization is empty.

\section{Transport of correspondences}
\label{ed:sec:transport}

The determinant calculation shows that $E_X$ becomes invertible over
$R^{\mathrm{loc}}=R[\zeta^{-1}]$. Indeed,
$e_\zeta=(-\zeta)^d(1+N)$ with $N$ in positive cohomological degree.
$N$ is nilpotent, so $(1+N)^{-1}$ is a finite sum. The factor $U_X$
is already invertible over $R$.

For smooth projective complex varieties $X,Y$, a cohomology kernel
$K_{YX}\in H^*(Y\times X;\mathbb Q)\otimes R$ defines
\[
 T_K(a)=(p_Y)_*(K_{YX}\smile p_X^*a).
\]
Pushforward uses the complex orientation. We compose kernels in the
same order as operators: on $Z\times Y\times X$,
\[
 K_{ZY}*K_{YX}
 =(p_{ZX})_*(p_{ZY}^*K_{ZY}\smile p_{YX}^*K_{YX}).
\]
All projections used for pushforward are proper. These formulas extend
$R^{\mathrm{loc}}$-linearly, and the diagonal class is the identity kernel.

\begin{proposition}\label{ed:prop:transport}
The transformed kernel
\begin{equation}\label{ed:eq:kernel}
 \widetilde K_{YX}
 =p_Y^*\epsilon_Y\smile K_{YX}\smile p_X^*\epsilon_X^{-1}
\end{equation}
acts by $E_YT_KE_X^{-1}$. This transformation preserves convolution,
identity kernels, and every operator relation over $R^{\mathrm{loc}}$.
If the original operators are defined over $R$, their transforms act
over $R$ on the image modules
\begin{equation}\label{ed:eq:image-module}
 M_X=E_X(H_{X,R})\subset H_{X,R}.
\end{equation}
\end{proposition}

\begin{proof}
The projection formula gives the operator identity. In a convolution,
the two factors pulled back from the middle variety are
$\epsilon_Y^{-1}\epsilon_Y=1$. They are even, so no sign is introduced.
This proves preservation of convolution. Restriction to the diagonal
similarly cancels the two factors and preserves its class.

Finally, injectivity of $E_X$ identifies the free module $H_{X,R}$
with its image $M_X$. Through this isomorphism the transformed action
is simply $E_YT_KE_X^{-1}:M_X\to M_Y$, which is defined over $R$.
\end{proof}

The image module is free, although its quotient in the ambient module
has $\zeta$-torsion. Regularity on the image therefore does not imply
regularity on the ambient module. Moreover, multiplication by
$\epsilon_X$ is a module map. It is generally not a unital algebra map
for the original cup product. The proposition is an operator transport
statement, with no multiplicative identification of coefficient algebras.

\section{Simultaneous preservation and the vacuum}
\label{ed:sec:order}

Let $L_0=R_0^D$ and let $M\subset L_0$ be a submodule of full rank.
Choose a Smith basis so that
$M=D_\ell L_0$, where
$D_\ell=\operatorname{diag}(\zeta^{\ell_1},\ldots,\zeta^{\ell_D})$.
The algebra of endomorphisms preserving both modules is
\[
 \mathscr O_\ell=
 \{T\in\operatorname{End}_{R_0}(L_0):T(M)\subset M\}.
\]

\begin{proposition}\label{ed:prop:order}
An ambient endomorphism belongs to $\mathscr O_\ell$ precisely when
\begin{equation}\label{ed:eq:order}
 T_{ij}\in\zeta^{\max(0,\ell_i-\ell_j)}R_0.
\end{equation}
For $\ell=(0,m)$ with $m>0$, the algebra
$\mathscr O_\ell/\zeta\mathscr O_\ell$ has dimension four. It has two
orthogonal idempotents and two opposite arrows, with both arrow
composites zero. Its image acting on $L_0/\zeta L_0$ has dimension three.
\end{proposition}

\begin{proof}
Preservation of $M$ is equivalent to regularity of
$D_\ell^{-1}TD_\ell$. Its $(i,j)$ entry is
$\zeta^{\ell_j-\ell_i}T_{ij}$. Combining this with ambient regularity
gives \eqref{ed:eq:order}.

For two exponents the algebra is
\[
 \begin{pmatrix}R_0&R_0\\\zeta^mR_0&R_0\end{pmatrix}
\]
with free basis $e_{11},e_{22},e_{12},\zeta^me_{21}$.
The off-diagonal products are $\zeta^me_{11}$ and
$\zeta^me_{22}$. Both vanish modulo $\zeta\mathscr O_\ell$.
The lower arrow has nonzero class in this abstract fiber, but acts as
zero on $L_0/\zeta L_0$. The other three basis elements act independently.
\end{proof}

Applied to the Euler image, this condition determines which regular
ambient operators descend to the projective-bundle quotient. The unit
factorization makes the same calculation valid after extension to $R$.
The distinction between the abstract fiber and its action is essential:
a generator may survive in the algebra while its specialized operator
vanishes.

For a K3 or Enriques surface $Y$, put $X_n=Y^{[n]}$ and
$H_n=H^*(X_n;\mathbb Q)\otimes R$, with empty vacuum
$|0\rangle=1\in H_0=R$. Normalize the length-one Nakajima creator and
annihilator so that
\[
 C_1(\omega)|0\rangle=\omega,\qquad
 [A_1(1),C_1(\omega)]|0\rangle=|0\rangle.
\]
The annihilator lowers length and kills the empty vacuum. Here both
insertions are even, so the bracket is the ordinary commutator.
Let $\widetilde C_1,\widetilde A_1$ denote their transports from
Proposition~\ref{ed:prop:transport}.

\begin{proposition}\label{ed:prop:vacuum}
The transported creator and the specified empty vacuum cannot be part
of a regular Heisenberg action on the unchanged modules $H_n$ at
$\zeta=0$. The transported original correspondence action is regular
on the image modules $M_n=E_{X_n}(H_n)$.
\end{proposition}

\begin{proof}
Since $E_{X_0}=1$, the transported creator sends the vacuum to
\[
 \widetilde C_1(\omega)|0\rangle
 =\epsilon_Y\omega=\zeta^2u(q,\zeta)\omega,
 \qquad u(q,\zeta)\in R^\times.
\]
Only the degree-zero component of the unit $U_Y$ contributes to this
product, and $c_2(T_Y)\omega=0$.
Any $R$-linear annihilator $H_1\to H_0$ sends this vector into
$\zeta^2R$. The reversed composition on the vacuum is zero, because
there is no negative length. The normalized Heisenberg relation would
make their difference equal to $1\in R$, a contradiction.
Regularity on $M_n$ follows from Proposition~\ref{ed:prop:transport}.
\end{proof}

These results determine the degeneration of the cohomological
correspondence representation. A compact critical Hall realization
requires an additional comparison with its coefficient complexes,
preserving the supported products, vacuum, orientations and reduction
maps. Likewise, the formal parameter used here does not construct the
modular line over the genus-two moduli stack. The Euler quotient
provides a concrete recipient for these comparisons and a constraint
on any comparison that retains the ambient cohomology and the vacuum.
